\section{Zadanie 11}
W urnie znajdują się dwie kule: jedna czarna i jedna biała. W każdej z \(n-2\) rund losujemy w sposób jednostajny jedną kulę z urny i dokładamy do urny kulę o tym samym kolorze co wylosowana. Pokaż, że liczba białych kul w urnie na końcu jest z jednakowym prawdopodobieństwem jedną z liczb ze zbioru \([n-1].\)

\subsection*{Rozwiązanie}
\(k\) – końcowa liczba białych kul w urnie.\\[1em]

Wszystkich możliwości losowania mamy \((n-1)!\) (najpierw losujemy z 2 kul, później z 3 itd. aż do \(n-1\)).\\[1em]

Policzymy teraz, ile jest możliwości losowań, aby końcowo liczba białych kul w urnie była równa \(k\).\\[0.5em]

\begin{itemize}
	\item \(\displaystyle \binom{n-2}{k-1}\) – wybór rund, w których wylosowano białe kule,\\
	\item \((k-1)!\) – możliwości wylosowania białych kul,\\
	\item \((n-k-1)!\) – możliwości wylosowania czarnych kul.
\end{itemize}

Zatem prawdopodobieństwo, że na koniec w urnie będzie \(k\) białych kul wynosi:
\[
	\frac{\binom{n-2}{k-1} (k-1)! (n-k-1)!}{(n-1)!}
	= \frac{(n-2)! (k-1)! (n-k-1)!}{(k-1)! (n-k-1)! (n-1)!}
	= \frac{1}{n-1}.
\]

Niezależnie od wartości \(k\), prawdopodobieństwo wynosi \(\dfrac{1}{n-1}\).
